% Compile with: uplatex file.tex && dvipdfmx file.dvi
%           or: latexmk -uplatex file.tex
%\documentclass[uplatex,dvipdfmx]{jsarticle}
\documentclass[dvipdfmx,a4paper,11pt]{jsarticle}
\usepackage{amsmath, amssymb}
\usepackage{tikz}

\title{scenario\_information\_theory}
\date{\today}

\begin{document}
\maketitle

シャノンエントロピー $H=-\sum_x p(x)\log_2 p(x)$ を、公平なサイコロ（$p=1/6$、6面）について定義そのままの総和で bit 単位で書く。専用操作は使わない。

$H = -\sum_{x=1}^{6} \frac{1}{6}\log_2 \frac{1}{6}$

\begin{equation}
  H = -\sum_{x=1}^{6} \frac{1}{6}\log_2 \frac{1}{6}
  \label{eq:E1}
\end{equation}

右辺を評価すると $\log 6/\log 2 = \log_2 6 \approx 2.585$ bit。情報理論の既定の単位は bit なので、これが本来の姿である。底つきの綴りは 2026-08-16 まで `log(2)*6`（2の対数$\times$6）という silent-wrong だった。

E1 の右辺を計算して

\begin{equation}
  -\sum_{x=1}^{6} \frac{1}{6}\log_2 \frac{1}{6} = \frac{\log 6}{\log 2}
  \label{eq:E2}
\end{equation}

次にベルヌーイ分布のエントロピー $H_b(p)$ を書く。ここからは一般操作が自然につながる。

$H_b = -p\log p - (1-p)\log(1-p)$

\begin{equation}
  H_b = -p\log p - (1-p)\log(1-p)
  \label{eq:E3}
\end{equation}

$p$ で微分して、エントロピーが最大になる条件を探す。

E3 を p で微分して

\begin{equation}
  \frac{dH_b}{dp} = \log \left(1-p\right)-\log p
  \label{eq:E4}
\end{equation}

導関数を 0 とおいて $p$ について解く。CAS は $\log p=\log(1-p)$ までを返す——対数の単射性を使う一歩は踏まないので、$p=1/2$ は研究者が読む。補わないのが正しい。

E4 を p について解いて

\begin{equation}
  \log p = \log \left(1-p\right)
  \label{eq:E5}
\end{equation}

Kullback–Leibler ダイバージェンスを、2値分布の間で定義そのままに書く。被和項の $\log\frac{p}{q}$ がそのまま書けることが要点。

$D = p\log \frac{p}{q} + (1-p)\log \frac{1-p}{1-q}$

\begin{equation}
  D = p\log \frac{p}{q} + (1-p)\log \frac{1-p}{1-q}
  \label{eq:E6}
\end{equation}

$q$ に $p$ を代入する。$D(p\Vert p)=0$ という KL の基本性質が、定義の表示ではなく計算として出る。

E6 の q に p を代入して

\begin{equation}
  D = 0
  \label{eq:E7}
\end{equation}

相互情報量を二重和の定義そのままに書く。

$I = \sum_{x=1}^{2}\sum_{y=1}^{2} p_{xy}(x,y)\log \frac{p_{xy}(x,y)}{p_x(x) p_y(y)}$

\begin{equation}
  I = \sum_{x=1}^{2}\sum_{y=1}^{2} p_{xy}(x,y)\log \frac{p_{xy}(x,y)}{p_x(x) p_y(y)}
  \label{eq:E8}
\end{equation}

独立な場合、すなわち $p_{xy}(x,y)=p_x(x)p_y(y)$ を入れた形を書く。

$I_0 = \sum_{x=1}^{2}\sum_{y=1}^{2} p_x(x) p_y(y)\log \frac{p_x(x) p_y(y)}{p_x(x) p_y(y)}$

\begin{equation}
  I_0 = \sum_{x=1}^{2}\sum_{y=1}^{2} p_x(x) p_y(y)\log \frac{p_x(x) p_y(y)}{p_x(x) p_y(y)}
  \label{eq:E9}
\end{equation}

右辺を評価すると 0。独立なら相互情報量が消えることが計算で確かめられる。

E9 の右辺を計算して

\begin{equation}
  \sum_{x=1}^{2}\sum_{y=1}^{2} p_x(x) p_y(y)\log \frac{p_x(x) p_y(y)}{p_x(x) p_y(y)} = 0
  \label{eq:E10}
\end{equation}

\end{document}